\documentclass{article}
\usepackage[margin=0.8in]{geometry}
\usepackage{amsmath,amssymb,amsthm,mathtools}
\usepackage{mathrsfs,bm,bbm}
\usepackage{graphicx,booktabs,array,multirow,tabularx,longtable}
\usepackage{enumitem}
\usepackage{xcolor}
\usepackage{algorithm,algpseudocode}
\usepackage{subcaption,tikz}
\usepackage{siunitx,cancel,accents}
\usepackage{microtype}
\usepackage[numbers,sort&compress]{natbib}
\usepackage[colorlinks=true,linkcolor=blue,citecolor=blue,urlcolor=blue]{hyperref}
\usepackage{cleveref}
\setlength{\emergencystretch}{2em}
\newtheorem{assumption}{Assumption}
\newtheorem{definition}{Definition}
\newtheorem{theorem}{Theorem}
\newtheorem{lemma}{Lemma}
\newtheorem{proposition}{Proposition}
\newtheorem{openendedquestion}{Open-ended Question}
\newtheorem{conjecture}{Conjecture}
\newcommand{\coreref}[1]{\ref{#1}}

\begin{document}

\sloppy

\section*{pid\_\allowbreak{}pareto\_\allowbreak{}benefit\_\allowbreak{}capacity}

\noindent\textit{Specialization:} v1\par

\paragraph{TL;DR.} A complete trial roster and finite positive mean roster size identify individual-weighted vector response-pattern laws on \(\mathcal C_{+}=\{c:\kappa(c)>0\}\); null strata contribute no mass and enter neither conditional laws nor linear-program rows. The overlapping projection equations define nonrectangular feasible arm-state distributions, which must be optimized jointly with the same-person cross-world coupling. The central theorem gives the attained sharp Pareto-benefit interval and joint dual certificates for every finite compatible sensitivity index. More generally, any prespecified finite family of cross-world payoffs has an attained sharp joint region equal to the linear image of \(\mathcal F_{\Gamma}(H)\), a compact convex polytope with an attained LP dual in every direction. A distinct breakdown point \(\Gamma_{\mathrm{maj}}\) marks when a compatible law first permits at most one-half of patients to benefit, and a three-binary-outcome internal diagnostic has joint upper endpoint \(8/9\) but rectangular-relaxation endpoint \(1\).

\section{Project justification}

\paragraph{Gap.} Cluster-trial missing-data methods target smooth average effects; Scharfstein, Manski, and Anthony bound scalar ordinal follow-up distributions; and randomized-trial fraction-benefiting methods treat scalar ordinal or binary outcomes. Gao, Ge, and Qian use fixed identified arm marginals. Fan, Park, Pass, and Shi use two identified source laws \((Y_1,X)\) and \((Y_0,X)\). Neither optimal-transport information structure preserves the nonrectangular restrictions from overlapping vector response patterns while simultaneously completing latent arm-state laws and ranging over the same-person potential-outcome coupling.

\paragraph{Niche.} Finite multivariate outcomes with nonmonotone component response patterns create a tractable but genuinely joint observed-data and cross-world partial-identification problem.

\paragraph{Fill.} The project gives the exact feasible polytope on positive-mass response strata, its compatibility onset, an attained sharp Pareto-benefit interval with primal observational embeddings and joint dual certificates, and the attained compact-polytope image for every prespecified finite family of cross-world linear payoffs with a dual certificate in each direction. It also gives an internal finite diagnostic proving strict loss from statewise rectangularization and a separate response-selection breakdown threshold for guaranteed majority benefit.

\section{Related work}

Wang, Li, and Wang provide the restricted-missing-at-random benchmark for incomplete outcomes and informative cluster sizes, but their estimand is a smooth average effect; their work informs the roster and missingness contract. Scharfstein, Manski, and Anthony study an ordinal outcome measured twice with the second measurement missing: Sections 2--4, especially equations (9)--(10) and (12)--(20), bound the scalar follow-up law under monotonicity and indexed selection restrictions; Section 5.2 discusses infeasible empirical constraint systems; and Section 6, equations (23)--(24), compares interventions through an arm-distribution stochastic-dominance measure. That program neither has overlapping vector response patterns nor couples two potential outcome vectors for a same-person strict-Pareto event. Huang, Fang, Hanley, and Rosenblum construct confidence intervals for a scalar ordinal or binary fraction benefiting in randomized trials; their target differs from vector strict-Pareto benefit under response-pattern selection and a jointly optimized missing-state and cross-world set. Gao, Ge, and Qian define the causal set in Section 1.2, equation (2), by ranging over couplings whose \(K\) arm marginals are fixed, and Section 2.2, equation (3), retains those fixed-marginal constraints for multi-marginal optimal transport. Their exact fixed-marginal MOT cannot encode arm marginals that are themselves endogenous completion variables tied across states by overlapping response-pattern projection equations. Fan, Park, Pass, and Shi study a different incomplete-data structure: Section 1.1, equation (1), defines a moment model in \((Y_1,Y_0,X)\); Section 2, Assumption 2.1(ii)--(iii), takes the laws of \((Y_1,X)\) and \((Y_0,X)\) as identified from two data sources; and Theorem 3.1 characterizes the resulting set by conditional optimal transport. The models are non-nested: their source laws are identified but no within-vector response-pattern completion is present, whereas the present response-pattern projections only partially identify each latent complete arm state before the cross-world coupling is chosen. Thus the present proposal does not claim that generic incomplete-data conditional optimal transport is absent. Ji, Lei, and Spector and Ober-Reynolds likewise inform the coupling component after arm laws or conditional arm laws are available. Strassen supplies the classical marginal-coupling principle behind the complete-response check. Li, Miao, Shpitser, and Tchetgen Tchetgen identify multivariate nonmonotone missing-data laws under self-censoring and completeness, a different point-identification model. The rectangular-relaxation construction is solely an internal diagnostic and is not attributed to a published imputation procedure. No theorem-level improvement over these comparator results is claimed.

\section{Setup}

Target (estimand / structural parameter): \(\beta=\mathbb P^{\star}\{Y_{ij}(1)\text{ strictly Pareto-dominates }Y_{ij}(0)\}\).
Primary functional / estimator / diagnostic: \(\mathcal I_{\Gamma}(H)=[\min_{x\in\mathcal F_{\Gamma}(H)}c_R^{\mathsf T}x,\ \max_{x\in\mathcal F_{\Gamma}(H)}c_R^{\mathsf T}x],\quad \Gamma_{\mathrm{feas}}\le\Gamma<\infty\).
\begin{itemize}
  \item \texttt{d} --- \texttt{positive integer}; \(\mathbb N\)
  \par\smallskip\noindent\emph{Definition.} \texttt{The prespecified number of outcome coordinates.}
  \item \texttt{k} --- \texttt{outcome-\allowbreak{}coordinate index}; \(\{1,\ldots,d\}\)
  \item \texttt{G} --- \texttt{positive integer}; \(\mathbb N\)
  \par\smallskip\noindent\emph{Definition.} \texttt{The number of randomized clusters.}
  \item \texttt{i} --- \texttt{cluster index}; \(\{1,\ldots,G\}\)
  \item \texttt{N\_\allowbreak{}i} --- \texttt{roster size}; \(\mathbb N\); Number of labeled individuals in cluster \(i\).
  \item \texttt{j} --- \texttt{individual index}; \(\{1,\ldots,N_i\}\)
  \item \texttt{\textbackslash{}\allowbreak{}mathcal R\_\allowbreak{}i} --- \texttt{complete labeled preassignment roster}
  \par\smallskip\noindent\emph{Definition.} The labels and all preassignment variables for the \(N_i\) target individuals in cluster \(i\).
  \item \texttt{\textbackslash{}\allowbreak{}mathcal B} --- \texttt{measurable baseline-\allowbreak{}record space}
  \par\smallskip\noindent\emph{Definition.} \texttt{The support of complete cluster-\allowbreak{}level preassignment records,\allowbreak{} including variable-\allowbreak{}length rosters.}
  \item \texttt{b} --- \texttt{baseline-\allowbreak{}record value}; \(\mathcal B\)
  \item \texttt{B\_\allowbreak{}i} --- \texttt{baseline record}; \(\mathcal B\)
  \par\smallskip\noindent\emph{Definition.} The observed preassignment record containing \(N_i\), \(\mathcal R_i\), and every roster variable used by assignment or response models.
  \item \texttt{\textbackslash{}\allowbreak{}mathcal C} --- \texttt{finite set}
  \par\smallskip\noindent\emph{Definition.} \texttt{The prespecified finite support of response-\allowbreak{}model strata.}
  \item \texttt{c} --- \texttt{response-\allowbreak{}stratum value}; \(\mathcal C\)
  \item \texttt{g\_\allowbreak{}\{\textbackslash{}\allowbreak{}mathrm\{resp\}\allowbreak{}\}\allowbreak{}} --- \texttt{prespecified response-\allowbreak{}stratification map}; \(g_{\mathrm{resp}}:\mathcal B\to\mathcal C\)
  \item \texttt{C\_\allowbreak{}i} --- \texttt{response-\allowbreak{}model stratum}; \(\mathcal C\)
  \par\smallskip\noindent\emph{Definition.} \(C_i=g_{\mathrm{resp}}(B_i)\).
  \item \texttt{A\_\allowbreak{}i} --- \texttt{cluster treatment}; \(\{0,1\}\)
  \item \texttt{a} --- \texttt{treatment-\allowbreak{}arm index}; \(\{0,1\}\)
  \item \texttt{e} --- \texttt{known assignment function}; \(e:\mathcal B\to(0,1)\)
  \item \texttt{e\_\allowbreak{}a} --- \texttt{arm assignment function}; \(e_a:\mathcal B\to(0,1)\)
  \par\smallskip\noindent\emph{Definition.} \(e_1(b)=e(b)\) and \(e_0(b)=1-e(b)\).
  \item \texttt{\textbackslash{}\allowbreak{}mathcal S\_\allowbreak{}k} --- \texttt{finite ordered coordinate space}
  \par\smallskip\noindent\emph{Definition.} The prespecified finite support of outcome coordinate \(k\), oriented so that larger values are better.
  \item \texttt{\textbackslash{}\allowbreak{}mathcal S} --- \texttt{finite product-\allowbreak{}ordered outcome space}; \(\mathcal S=\prod_{k=1}^{d}\mathcal S_k\)
  \par\smallskip\noindent\emph{Definition.} \texttt{Each finite coordinate set is oriented so that larger values are better.}
  \item \texttt{s} --- \texttt{complete outcome state}; \(\mathcal S\)
  \item \texttt{s\_\allowbreak{}0} --- \texttt{control outcome state}; \(\mathcal S\)
  \item \texttt{s\_\allowbreak{}1} --- \texttt{treated outcome state}; \(\mathcal S\)
  \item \texttt{Y\_\allowbreak{}\{ij\}\allowbreak{}(a)} --- \texttt{vector potential outcome}; \(\mathcal S\)
  \item \texttt{\textbackslash{}\allowbreak{}mathcal M} --- \texttt{response-\allowbreak{}pattern space}; \(\mathcal M=\{0,1\}^{d}\)
  \item \texttt{m} --- \texttt{response-\allowbreak{}pattern value}; \(\mathcal M\)
  \item \texttt{\textbackslash{}\allowbreak{}mathcal S\_\allowbreak{}m} --- \texttt{finite projected-\allowbreak{}state space}; \(\mathcal S_m=\prod_{k:m_k=1}\mathcal S_k\)
  \item \texttt{M\_\allowbreak{}\{ij\}\allowbreak{}(a)} --- \texttt{potential vector response pattern}; \(\mathcal M\)
  \item \texttt{\textbackslash{}\allowbreak{}operatorname\{pr\}\allowbreak{}\_\allowbreak{}m} --- \texttt{coordinate projection}; \(\operatorname{pr}_m:\mathcal S\to\mathcal S_m\)
  \par\smallskip\noindent\emph{Definition.} The projection retaining exactly the coordinates marked observed by \(m\), with \(\mathcal S_m=\prod_{k:m_k=1}\mathcal S_k\).
  \item \texttt{z\_\allowbreak{}m} --- \texttt{observed projected-\allowbreak{}state value}; \(\mathcal S_m\)
  \item \texttt{O\_\allowbreak{}\{ij\}\allowbreak{}} --- \texttt{observed individual record}
  \par\smallskip\noindent\emph{Definition.} \(O_{ij}=(B_i,A_i,M_{ij}(A_i),\operatorname{pr}_{M_{ij}(A_i)}Y_{ij}(A_i))\).
  \item \texttt{\textbackslash{}\allowbreak{}nu} --- \texttt{positive real}
  \par\smallskip\noindent\emph{Definition.} \(\nu=\mathbb E[N_i]\).
  \item \texttt{\textbackslash{}\allowbreak{}mathbb P\textasciicircum{}\{\textbackslash{}\allowbreak{}star\}\allowbreak{}} --- \texttt{individual-\allowbreak{}weighted probability law}
  \par\smallskip\noindent\emph{Definition.} The law obtained by sampling a cluster with probability proportional to \(N_i\) and then sampling one of its \(N_i\) roster members uniformly.
  \item \texttt{\textbackslash{}\allowbreak{}kappa} --- \texttt{individual-\allowbreak{}weighted response-\allowbreak{}stratum mass function}; \(\kappa:\mathcal C\to[0,1]\)
  \par\smallskip\noindent\emph{Definition.} \(\kappa(c)=\mathbb P^{\star}(C_i=c)\).
  \item \texttt{\textbackslash{}\allowbreak{}mathcal C\_\allowbreak{}\{+\allowbreak{}\}\allowbreak{}} --- \texttt{positive-\allowbreak{}mass response-\allowbreak{}stratum index set}; \(\mathcal C_{+}\subseteq\mathcal C\)
  \par\smallskip\noindent\emph{Definition.} \(\mathcal C_{+}=\{c\in\mathcal C:\kappa(c)>0\}\). Strata outside \(\mathcal C_{+}\) have zero individual-weighted mass and are omitted from every conditional-law and linear-program index set.
  \item \texttt{H\_\allowbreak{}a} --- \texttt{identified arm-\allowbreak{}specific response-\allowbreak{}pattern law}; \(H_a:\mathcal C_{+}\times\mathcal M\times\mathcal S_m\to[0,1]\); Observed-law right-hand side for arm \(a\) on positive-mass strata.
  \item \texttt{r\_\allowbreak{}a} --- \texttt{identified response-\allowbreak{}pattern probability}; \(r_a:\mathcal M\times\mathcal C_{+}\to[0,1]\); Conditional response-pattern frequency under arm \(a\).
  \item \texttt{\textbackslash{}\allowbreak{}Gamma} --- \texttt{sensitivity index}; \([1,\infty)\)
  \par\smallskip\noindent\emph{Definition.} \texttt{The maximal multiplicative departure of a response-\allowbreak{}pattern probability from its response-\allowbreak{}stratum marginal frequency.}
  \item \texttt{\textbackslash{}\allowbreak{}gamma} --- \texttt{candidate sensitivity index}; \([1,\infty)\)
  \item \texttt{\textbackslash{}\allowbreak{}lambda\_\allowbreak{}a} --- \texttt{latent full-\allowbreak{}state and response-\allowbreak{}pattern mass}; \(\lambda_a:\mathcal C_{+}\times\mathcal S\times\mathcal M\to[0,1]\); \texttt{Arm-\allowbreak{}specific completion variable.}
  \item \texttt{q\_\allowbreak{}a} --- \texttt{latent complete arm-\allowbreak{}state mass}; \(q_a:\mathcal C_{+}\times\mathcal S\to[0,1]\); \texttt{Individual-\allowbreak{}weighted joint mass of response stratum and complete potential state.}
  \item \texttt{\textbackslash{}\allowbreak{}pi} --- \texttt{response-\allowbreak{}stratum-\allowbreak{}preserving cross-\allowbreak{}world coupling}; \(\pi:\mathcal C_{+}\times\mathcal S\times\mathcal S\to[0,1]\); Joint individual-weighted law of \(C_i,Y_{ij}(0),Y_{ij}(1)\).
  \item \texttt{R} --- \texttt{strict Pareto-\allowbreak{}benefit relation}; \(R\subseteq\mathcal S\times\mathcal S\)
  \par\smallskip\noindent\emph{Definition.} \(R=\{(s_0,s_1):s_{1k}\ge s_{0k}\text{ for every }k,\ s_1\ne s_0\}\).
  \item \texttt{\textbackslash{}\allowbreak{}beta} --- \texttt{causal probability}; \([0,1]\)
  \par\smallskip\noindent\emph{Definition.} \(\beta=\mathbb P^{\star}((Y_{ij}(0),Y_{ij}(1))\in R)\).
  \item \texttt{x} --- \texttt{finite nonnegative decision vector}
  \par\smallskip\noindent\emph{Definition.} The vector concatenating \(\lambda_0,\lambda_1,q_0,q_1,\pi\) in a fixed lexicographic order.
  \item \texttt{\textbackslash{}\allowbreak{}bar h} --- \texttt{finite equality right-\allowbreak{}hand-\allowbreak{}side vector}
  \par\smallskip\noindent\emph{Definition.} The vector stacking \(H_0,H_1\) and zeros for all completion and coupling consistency equations.
  \item \texttt{E} --- \texttt{finite equality-\allowbreak{}constraint matrix}; Matrix encoding observed projections, arm-state completion, and coupling marginals on \(\mathcal C_{+}\).
  \item \texttt{G\_\allowbreak{}\{\textbackslash{}\allowbreak{}Gamma\}\allowbreak{}} --- \texttt{finite inequality-\allowbreak{}constraint matrix}; Matrix encoding the two-sided response-pattern sensitivity restrictions on \(\mathcal C_{+}\).
  \item \texttt{\textbackslash{}\allowbreak{}alpha} --- \texttt{unrestricted equality-\allowbreak{}dual vector}; \texttt{a finite-\allowbreak{}dimensional Euclidean space}
  \item \texttt{\textbackslash{}\allowbreak{}eta} --- \texttt{nonnegative sensitivity-\allowbreak{}dual vector}; \texttt{a finite-\allowbreak{}dimensional nonnegative orthant}
  \item \texttt{c\_\allowbreak{}R} --- \texttt{finite objective vector}
  \par\smallskip\noindent\emph{Definition.} The entries corresponding to \(\pi(c,s_0,s_1)\) equal \(\mathbf 1\{(s_0,s_1)\in R\}\), and all other entries equal zero.
  \item \texttt{\textbackslash{}\allowbreak{}mathcal F\_\allowbreak{}\{\textbackslash{}\allowbreak{}Gamma\}\allowbreak{}(H)} --- \texttt{joint observed-\allowbreak{}data feasible polytope}; All arm completions and cross-world couplings on \(\mathcal C_{+}\) compatible with the response-pattern law and sensitivity index.
  \item \texttt{\textbackslash{}\allowbreak{}Gamma\_\allowbreak{}\{\textbackslash{}\allowbreak{}mathrm\{feas\}\allowbreak{}\}\allowbreak{}} --- \texttt{compatibility onset}; \([1,\infty]\); \texttt{Smallest sensitivity index at which the observed response-\allowbreak{}pattern law admits a full-\allowbreak{}data completion and cross-\allowbreak{}world coupling.}
  \item \texttt{L\_\allowbreak{}\{\textbackslash{}\allowbreak{}Gamma\}\allowbreak{}} --- \texttt{sharp lower endpoint}; \([0,1]\); \texttt{Minimum Pareto-\allowbreak{}benefit mass over the joint feasible polytope.}
  \item \texttt{U\_\allowbreak{}\{\textbackslash{}\allowbreak{}Gamma\}\allowbreak{}} --- \texttt{sharp upper endpoint}; \([0,1]\); \texttt{Maximum Pareto-\allowbreak{}benefit mass over the joint feasible polytope.}
  \item \texttt{\textbackslash{}\allowbreak{}mathcal I\_\allowbreak{}\{\textbackslash{}\allowbreak{}Gamma\}\allowbreak{}(H)} --- \texttt{identified set}; closed subinterval of \([0,1]\); Observed-data identified set for \(\beta\).
  \item \texttt{\textbackslash{}\allowbreak{}ell\_\allowbreak{}a} --- \texttt{statewise lower envelope}; \(\ell_a:\mathcal C_{+}\times\mathcal S\to[0,1]\); Coordinatewise minimum of \(q_a(c,s)\) over arm-specific feasible completions.
  \item \texttt{u\_\allowbreak{}a} --- \texttt{statewise upper envelope}; \(u_a:\mathcal C_{+}\times\mathcal S\to[0,1]\); Coordinatewise maximum of \(q_a(c,s)\) over arm-specific feasible completions.
  \item \texttt{\textbackslash{}\allowbreak{}mathcal F\textasciicircum{}\{\textbackslash{}\allowbreak{}mathrm\{rect\}\allowbreak{}\}\allowbreak{}\_\allowbreak{}\{\textbackslash{}\allowbreak{}Gamma\}\allowbreak{}(H)} --- \texttt{rectangular completion relaxation}; \texttt{The relaxation retaining only statewise envelopes,\allowbreak{} normalization,\allowbreak{} and generic cross-\allowbreak{}world transport.}
  \item \texttt{\textbackslash{}\allowbreak{}mathsf W} --- \texttt{finite response-\allowbreak{}pattern witness}; \texttt{A three-\allowbreak{}binary-\allowbreak{}outcome law with pairwise observation patterns showing strict nonrectangularity.}
  \item \texttt{\textbackslash{}\allowbreak{}Gamma\_\allowbreak{}\{\textbackslash{}\allowbreak{}mathrm\{maj\}\allowbreak{}\}\allowbreak{}} --- \texttt{majority-\allowbreak{}benefit breakdown point}; \([1,\infty]\); \texttt{Smallest sensitivity index admitting a compatible law with at most one-\allowbreak{}half Pareto benefit.}
  \item \texttt{K} --- \texttt{positive integer}; \(\mathbb N\)
  \par\smallskip\noindent\emph{Definition.} \texttt{The number of prespecified payoff arrays in the joint-\allowbreak{}payoff corollary.}
  \item \texttt{h} --- \texttt{payoff index}; \(\{1,\ldots,K\}\)
  \item \texttt{w\_\allowbreak{}h} --- \texttt{real payoff array}; \(w_h:\mathcal C_{+}\times\mathcal S\times\mathcal S\to\mathbb R\); The \(h\)-th prespecified payoff on a response stratum and a pair of potential outcome states.
  \item \texttt{v} --- \texttt{direction vector}; \(\mathbb R^K\); \texttt{A direction in which the joint payoff identified polytope is supported.}
  \item \texttt{c\_\allowbreak{}v} --- \texttt{finite directional objective vector}
  \par\smallskip\noindent\emph{Definition.} For fixed \(v\) and \(w_1,\ldots,w_K\), the entry corresponding to \(\pi(c,s_0,s_1)\) is \(\sum_{h=1}^K v_h w_h(c,s_0,s_1)\), and every other entry is zero.
\end{itemize}

\section{Assumptions}

\begin{assumption}[ass:iid-clusters]\label{ass:iid-clusters}
The cluster records \(\{(B_i,\{Y_{ij}(a),M_{ij}(a):j\le N_i,a\in\{0,1\}\}) : i=1,\ldots,G\}\) are independent and identically distributed.
\par\smallskip\noindent\emph{Standard} (Independent superpopulation sampling of clusters, \cite{WangParkSmallLi2024}).
\end{assumption}

\begin{assumption}[ass:finite-mean-roster]\label{ass:finite-mean-roster}
\(0<\nu<\infty\).
\par\smallskip\noindent\emph{Standard} (Finite positive mean cluster size, \cite{WangLiWang2026}).
\end{assumption}

\begin{assumption}[ass:complete-roster]\label{ass:complete-roster}
The complete labeled roster \(\mathcal R_i\), its size \(N_i\), and every roster variable entering assignment or response models are observed before \(A_i\) and are included in \(B_i\).
\par\smallskip\noindent\emph{Novel.} A closed enrollment registry with prespecified baseline forms records every eligible label and every design or follow-up predictor before cluster assignment. This contract is empirically auditable and separates outcome nonresponse from unobserved recruitment.
\end{assumption}

\begin{assumption}[ass:cluster-randomization]\label{ass:cluster-randomization}
\(A_i\perp\!\!\perp\{Y_{ij}(a),M_{ij}(a):j\le N_i,a\in\{0,1\}\}\mid B_i\).
\par\smallskip\noindent\emph{Standard} (Conditional cluster randomization, \cite{WangParkSmallLi2024}).
\end{assumption}

\begin{assumption}[ass:known-assignment]\label{ass:known-assignment}
\(\mathbb P(A_i=a\mid B_i)=e_a(B_i)\) for \(a\in\{0,1\}\).
\par\smallskip\noindent\emph{Standard} (Known randomized assignment mechanism, \cite{WangParkSmallLi2024}).
\end{assumption}

\begin{assumption}[ass:assignment-positivity]\label{ass:assignment-positivity}
\(e_a(b)\ge\varepsilon\) for every \(a\in\{0,1\}\) and \(b\in\mathcal B\).
\par\smallskip\noindent\emph{Standard} (Treatment overlap, \cite{WangParkSmallLi2024}).
\end{assumption}

\begin{assumption}[ass:cluster-sutva]\label{ass:cluster-sutva}
Each \(Y_{ij}(a)\) and \(M_{ij}(a)\) depends only on its own cluster assignment \(a\).
\par\smallskip\noindent\emph{Standard} (Cluster-level stable unit treatment value assumption, \cite{HayesMoulton2017}).
\end{assumption}

\begin{assumption}[ass:observation-consistency]\label{ass:observation-consistency}
The realized response pattern and observed outcome components equal \(M_{ij}(A_i)\) and \(\operatorname{pr}_{M_{ij}(A_i)}Y_{ij}(A_i)\), respectively.
\par\smallskip\noindent\emph{Standard} (Potential-outcome consistency, \cite{Rubin1974}).
\end{assumption}

\begin{assumption}[ass:pattern-selection-sensitivity]\label{ass:pattern-selection-sensitivity}
For every \(a\in\{0,1\}\), \(c\in\mathcal C_{+}\), \(m\in\mathcal M\), and \(s\in\mathcal S\) with \(\mathbb P^{\star}(Y_{ij}(a)=s\mid C_i=c)>0\), \(\Gamma^{-1}r_a(m\mid c)\le\mathbb P^{\star}(M_{ij}(a)=m\mid C_i=c,Y_{ij}(a)=s)\le\Gamma r_a(m\mid c)\).
\par\smallskip\noindent\emph{Novel.} Follow-up audits can bound how much any complete clinical state is overrepresented or underrepresented in each vector response pattern among patients in the same prespecified response stratum. Because \(C_i=g_{\mathrm{resp}}(B_i)\), every stratifying variable and the full roster are recorded before assignment. The index has a direct likelihood-ratio interpretation, and \(\Gamma=1\) is vector restricted missing at random without requiring complete-case positivity.
\end{assumption}

\section{Definitions}

\begin{definition}[def:identified-response-laws]\label{def:identified-response-laws}
\par\smallskip\noindent\emph{Defined object.} \texttt{Identified response-\allowbreak{}pattern laws}
\par\smallskip\noindent\emph{Construction.} For \(c\in\mathcal C_{+}\), \(m\in\mathcal M\), and \(z_m\in\mathcal S_m\), set \(H_a(c,m,z_m)=\nu^{-1}\mathbb E[e_a(B_i)^{-1}\mathbf 1\{A_i=a,C_i=c\}\sum_{j=1}^{N_i}\mathbf 1\{M_{ij}(A_i)=m,\operatorname{pr}_mY_{ij}(A_i)=z_m\}]\) and \(r_a(m\mid c)=\kappa(c)^{-1}\sum_{z_m}H_a(c,m,z_m)\). No conditional response law is defined for \(c\notin\mathcal C_{+}\); such strata have zero \(\mathbb P^{\star}\)-mass and are omitted from the linear program.
\par\smallskip\noindent\emph{Inputs.} \texttt{ass:\allowbreak{}finite-\allowbreak{}mean-\allowbreak{}roster}; \texttt{ass:\allowbreak{}complete-\allowbreak{}roster}; \texttt{ass:\allowbreak{}cluster-\allowbreak{}randomization}; \texttt{ass:\allowbreak{}known-\allowbreak{}assignment}; \texttt{ass:\allowbreak{}assignment-\allowbreak{}positivity}; \texttt{ass:\allowbreak{}observation-\allowbreak{}consistency}.
\end{definition}

\begin{definition}[def:joint-feasible-polytope]\label{def:joint-feasible-polytope}
\par\smallskip\noindent\emph{Defined object.} \texttt{Joint missing-\allowbreak{}state and cross-\allowbreak{}world feasible polytope}
\par\smallskip\noindent\emph{Construction.} \(\mathcal F_{\Gamma}(H)\) is the set of nonnegative \(x=(\lambda_0,\lambda_1,q_0,q_1,\pi)\) satisfying, for every \(a\in\{0,1\}\), \(c\in\mathcal C_{+}\), \(m\in\mathcal M\), \(s,z_m,s_0,s_1\) in their declared finite state spaces, \(\sum_{s:\operatorname{pr}_m(s)=z_m}\lambda_a(c,s,m)=H_a(c,m,z_m)\), \(q_a(c,s)=\sum_m\lambda_a(c,s,m)\), \(\Gamma^{-1}r_a(m\mid c)q_a(c,s)\le\lambda_a(c,s,m)\le\Gamma r_a(m\mid c)q_a(c,s)\), \(\sum_{s_1}\pi(c,s_0,s_1)=q_0(c,s_0)\), and \(\sum_{s_0}\pi(c,s_0,s_1)=q_1(c,s_1)\). Equivalently, \(\mathcal F_{\Gamma}(H)=\{x\ge0:Ex=\bar h,\ G_{\Gamma}x\le0\}\), where \(E\) and \(G_{\Gamma}\) are the finite coefficient matrices of the displayed constraints indexed by \(\mathcal C_{+}\), \(\mathcal M\), and \(\mathcal S\). Strata in \(\mathcal C\setminus\mathcal C_{+}\) contribute zero individual-weighted mass and generate no completion, coupling, or sensitivity rows.
\par\smallskip\noindent\emph{Inputs.} \texttt{def:\allowbreak{}identified-\allowbreak{}response-\allowbreak{}laws}; \texttt{ass:\allowbreak{}pattern-\allowbreak{}selection-\allowbreak{}sensitivity}.
\end{definition}

\begin{definition}[def:feasibility-onset]\label{def:feasibility-onset}
\par\smallskip\noindent\emph{Defined object.} \texttt{Sensitivity-\allowbreak{}model compatibility onset}
\par\smallskip\noindent\emph{Construction.} \(\Gamma_{\mathrm{feas}}=\inf\{\gamma\in[1,\infty):\mathcal F_{\gamma}(H)\ne\varnothing\}\), with \(\inf\varnothing=\infty\).
\par\smallskip\noindent\emph{Inputs.} \texttt{def:\allowbreak{}joint-\allowbreak{}feasible-\allowbreak{}polytope}.
\end{definition}

\begin{definition}[def:sharp-endpoints]\label{def:sharp-endpoints}
\par\smallskip\noindent\emph{Defined object.} \texttt{Joint sharp endpoints}
\par\smallskip\noindent\emph{Construction.} For every finite \(\Gamma\ge\Gamma_{\mathrm{feas}}\), define \(L_{\Gamma}=\min_{x\in\mathcal F_{\Gamma}(H)}c_R^{\mathsf T}x\), \(U_{\Gamma}=\max_{x\in\mathcal F_{\Gamma}(H)}c_R^{\mathsf T}x\), and \(\mathcal I_{\Gamma}(H)=[L_{\Gamma},U_{\Gamma}]\).
\par\smallskip\noindent\emph{Inputs.} \texttt{def:\allowbreak{}joint-\allowbreak{}feasible-\allowbreak{}polytope}; \texttt{def:\allowbreak{}feasibility-\allowbreak{}onset}.
\end{definition}

\begin{definition}[def:rectangular-relaxation]\label{def:rectangular-relaxation}
\par\smallskip\noindent\emph{Defined object.} \texttt{Statewise-\allowbreak{}envelope rectangular relaxation}
\par\smallskip\noindent\emph{Construction.} For each arm and \(c\in\mathcal C_{+}\), let \(\ell_a(c,s)=\min q_a(c,s)\) and \(u_a(c,s)=\max q_a(c,s)\) over the arm-specific \((\lambda_a,q_a)\) block satisfying the projection, completion, and sensitivity constraints in \(\mathcal F_{\Gamma}(H)\). Define \(\mathcal F^{\mathrm{rect}}_{\Gamma}(H)\) as the set of nonnegative \((q_0,q_1,\pi)\) satisfying \(\ell_a(c,s)\le q_a(c,s)\le u_a(c,s)\), \(\sum_s q_a(c,s)=\kappa(c)\), \(\sum_{s_1}\pi(c,s_0,s_1)=q_0(c,s_0)\), and \(\sum_{s_0}\pi(c,s_0,s_1)=q_1(c,s_1)\) for every state index and every \(c\in\mathcal C_{+}\). This relaxation retains the statewise envelopes but discards the joint cross-state projection equalities.
\par\smallskip\noindent\emph{Inputs.} \texttt{def:\allowbreak{}joint-\allowbreak{}feasible-\allowbreak{}polytope}.
\end{definition}

\begin{definition}[def:nonrectangular-witness]\label{def:nonrectangular-witness}
\par\smallskip\noindent\emph{Defined object.} \texttt{Three-\allowbreak{}coordinate pair-\allowbreak{}pattern witness}
\par\smallskip\noindent\emph{Construction.} In one response stratum, let \(\mathcal S=\{0,1\}^3\) in lexicographic order and let the three response patterns \(110,101,011\) each have probability \(1/3\), independently of the full state. Conditional on these patterns, respectively, the control projection masses on \(00,01,10,11\) are \((1,1,2,2)/6\), \((1,1,3,1)/6\), and \((2,1,2,1)/6\); the treatment projection masses are \((0,0,2,4)/6\), \((0,0,3,3)/6\), and \((1,1,2,2)/6\). The complete arm-state masses \(q_0=(1,0,0,1,1,1,2,0)/6\) and \(q_1=(0,0,0,0,1,1,2,2)/6\) generate these projections. Denote this observed law by \(\mathsf W\).
\par\smallskip\noindent\emph{Inputs.} \texttt{def:\allowbreak{}identified-\allowbreak{}response-\allowbreak{}laws}.
\end{definition}

\begin{definition}[def:majority-breakdown]\label{def:majority-breakdown}
\par\smallskip\noindent\emph{Defined object.} \texttt{Sensitivity threshold for guaranteed majority benefit}
\par\smallskip\noindent\emph{Construction.} \(\Gamma_{\mathrm{maj}}=\inf\{\gamma\in[1,\infty):\gamma\ge\Gamma_{\mathrm{feas}},\ \exists x\in\mathcal F_{\gamma}(H)\text{ with }c_R^{\mathsf T}x\le1/2\}\), with \(\inf\varnothing=\infty\).
\par\smallskip\noindent\emph{Inputs.} \texttt{def:\allowbreak{}joint-\allowbreak{}feasible-\allowbreak{}polytope}; \texttt{def:\allowbreak{}feasibility-\allowbreak{}onset}.
\end{definition}

\section{Main results}

\begin{proposition}[prop:feasibility-onset]\label{prop:feasibility-onset}
The finite-index feasible domain is empty exactly when \(\Gamma_{\mathrm{feas}}=\infty\). If \(\Gamma_{\mathrm{feas}}<\infty\), then the onset is attained and \(\mathcal F_{\Gamma}(H)\ne\varnothing\) holds exactly for finite \(\Gamma\ge\Gamma_{\mathrm{feas}}\). Thus \(\Gamma<\Gamma_{\mathrm{feas}}\) means that the indexed response-selection restriction is incompatible with the observed response-pattern law, rather than that the data fail to guarantee majority benefit.
\end{proposition}
\begin{proof}
Let
\[
\mathscr D=\{\gamma\in[1,\infty):\mathcal F_{\gamma}(H)\ne\varnothing\}.
\]
We first use monotonicity of the displayed sensitivity inequalities.  If
\(1\le\gamma\le\gamma'<\infty\) and \(x\in\mathcal F_{\gamma}(H)\), then, coordinate by coordinate,
\[
(\gamma')^{-1}r_a(m\mid c)q_a(c,s)
\le \gamma^{-1}r_a(m\mid c)q_a(c,s)
\le \lambda_a(c,s,m)
\le \gamma r_a(m\mid c)q_a(c,s)
\le \gamma' r_a(m\mid c)q_a(c,s).
\]
Here nonnegativity of \(r_a\) and \(q_a\) justifies both comparisons.  All equality constraints are independent of \(\gamma\), so \(x\in\mathcal F_{\gamma'}(H)\).  Thus \(\mathscr D\) is upward closed.

We next use compactness of finite probability tables.  Summing the projection equations over \(c,m,z_m\), and using the definition of the identified response laws together with known randomization, assignment positivity, and \(0<\nu<\infty\), gives
\[
\sum_{c,m,z_m}H_a(c,m,z_m)
=\nu^{-1}\mathbb E\!\left[e_a(B_i)^{-1}\mathbf 1\{A_i=a\}N_i\right]
=\nu^{-1}\mathbb E[N_i]=1.
\]
Consequently every feasible table satisfies
\[
\sum_{c,s,m}\lambda_a(c,s,m)=1,
\qquad
\sum_{c,s}q_a(c,s)=1,
\qquad
\sum_{c,s_0,s_1}\pi(c,s_0,s_1)=1.
\]
All coordinates of \(x\) therefore lie in \([0,1]\).  Because the state spaces are finite and all constraints are closed, every \(\mathcal F_{\gamma}(H)\) is compact.

Suppose \(\mathscr D\ne\varnothing\), and put \(\gamma_0=\inf\mathscr D\).  Since every element of \(\mathscr D\) is finite, \(1\le\gamma_0<\infty\).  Choose \(\gamma_n\in\mathscr D\) with \(\gamma_n\to\gamma_0\), and choose \(x_n\in\mathcal F_{\gamma_n}(H)\).  The preceding common bound gives a convergent subsequence, say \(x_n\to x_0\).  The equality rows pass to the limit.  Since \(\gamma_0\ge1\), both coefficient maps \(\gamma\mapsto\gamma\) and \(\gamma\mapsto\gamma^{-1}\) are continuous at \(\gamma_0\); hence every two-sided sensitivity inequality also passes to the limit.  Thus \(x_0\in\mathcal F_{\gamma_0}(H)\), so the onset is attained.  Upward closure then gives \(\mathscr D=[\gamma_0,\infty)\).  If \(\mathscr D=\varnothing\), its infimum is \(\infty\) by definition; conversely, any member of \(\mathscr D\) would make its infimum finite.  This proves all assertions, including that an index below the onset is an infeasible restriction rather than a causal conclusion.
\end{proof}
\par\smallskip\noindent \textit{Justification.} The onset separates model-data incompatibility from an adverse causal conclusion along the sensitivity path. \textit{Gap.} Scharfstein, Manski, and Anthony, Section 5.2, note that empirically constrained missing-ordinal programs can be infeasible and propose penalized slacks; the present population-law statement instead characterizes the exact no-slack compatibility domain of the joint response-pattern and cross-world program. \textit{Consumer.} A trial report can flag sensitivity values below the onset as refuted by the response-pattern law before interpreting any Pareto-benefit bound.

\begin{theorem}[thm:joint-observed-data-sharpness]\label{thm:joint-observed-data-sharpness}
For every finite \(\Gamma\ge\Gamma_{\mathrm{feas}}\), under the stated information contract, the sharp identified set for \(\beta\) over all compatible cluster potential-outcome and vector-response laws satisfying the \(\Gamma\)-selection restriction is exactly \(\mathcal I_{\Gamma}(H)=[L_{\Gamma},U_{\Gamma}]\). All conditional response laws and all rows of \(\mathcal F_{\Gamma}(H)\) are indexed only by \(c\in\mathcal C_{+}\); strata with \(\kappa(c)=0\) contribute zero individual-weighted mass and require no conditional-law version. Both endpoints and every intermediate value are attained. Every compatible law maps into \(\mathcal F_{\Gamma}(H)\), and every point of \(\mathcal F_{\Gamma}(H)\) can be embedded in a cluster law that reproduces \(H_0,H_1\) on \(\mathcal C_{+}\), the observed complete roster distribution, and the known randomization. Moreover, finite-dimensional strong duality gives attained certificates \(L_{\Gamma}=\max_{\alpha\text{ unrestricted},\ \eta\ge0}\{\bar h^{\mathsf T}\alpha:E^{\mathsf T}\alpha-G_{\Gamma}^{\mathsf T}\eta\le c_R\}\) and \(U_{\Gamma}=\min_{\alpha\text{ unrestricted},\ \eta\ge0}\{\bar h^{\mathsf T}\alpha:E^{\mathsf T}\alpha+G_{\Gamma}^{\mathsf T}\eta\ge c_R\}\).
\end{theorem}
\begin{proof}
Fix a finite \(\Gamma\geq\Gamma_{\mathrm{feas}}\).  Proposition \texttt{prop:feasibility-onset} implies that \(\mathcal F_{\Gamma}(H)\) is nonempty.  To distinguish the causal target from the observed-law functional, let \(P_{\mathrm{obs}}\) be the law of the observable cluster record
\[
O_i=\left(B_i,A_i,
  \left\{M_{ij}(A_i),
  \operatorname{pr}_{M_{ij}(A_i)}Y_{ij}(A_i):1\leq j\leq N_i\right\}
  \right).
\]
The information contract retains from \(P_{\mathrm{obs}}\) the distribution of \(B_i\), the known function \(e\), and the two finite arrays \(H_0,H_1\) in Definition \texttt{def:identified-response-laws}.  Let \(\Theta_I^{\Gamma}(H)\) be the set of values of the causal probability
\[
\mathbb P^{\star}\{(Y_{ij}(0),Y_{ij}(1))\in R\}
\]
over cluster potential-outcome and response laws that preserve those observed-law quantities and satisfy the assumptions listed in the theorem.  We prove that this set, rather than merely an observed regression functional, equals the claimed interval.

\emph{Technique: inverse-randomization identification and finite-table marginalization.}
Consider an arbitrary compatible law and define
\[
\begin{aligned}
\lambda_a(c,s,m)
  &=\mathbb P^{\star}\{C_i=c,Y_{ij}(a)=s,M_{ij}(a)=m\},\\
q_a(c,s)
  &=\mathbb P^{\star}\{C_i=c,Y_{ij}(a)=s\},\\
\pi(c,s_0,s_1)
  &=\mathbb P^{\star}\{C_i=c,Y_{ij}(0)=s_0,Y_{ij}(1)=s_1\}.
\end{aligned}
\]
These arrays are nonnegative.  Marginalization gives
\[
q_a(c,s)=\sum_{m\in\mathcal M}\lambda_a(c,s,m),
\quad
\sum_{s_1\in\mathcal S}\pi(c,s_0,s_1)=q_0(c,s_0),
\quad
\sum_{s_0\in\mathcal S}\pi(c,s_0,s_1)=q_1(c,s_1).
\]

For fixed \(a,c,m,z_m\), set
\[
V_i=\mathbf 1\{C_i=c\}\sum_{j=1}^{N_i}
 \mathbf 1\{M_{ij}(a)=m,\operatorname{pr}_mY_{ij}(a)=z_m\}.
\]
The complete-roster assumption makes \(C_i\) and the summation range preassignment quantities.  Conditional cluster randomization and the known assignment mechanism imply
\[
\mathbb E\!\left[
 \left.\frac{\mathbf 1\{A_i=a\}}{e_a(B_i)}V_i
 \right|B_i,\{Y_{ij}(u),M_{ij}(u):j\leq N_i,u\in\{0,1\}\}\right]=V_i.
\]
The division is valid because assignment positivity gives \(e_a(B_i)\geq\varepsilon>0\).  Moreover \(0\leq V_i\leq N_i\), so the expectation is finite by \(0<\nu=\mathbb E[N_i]<\infty\).  Cluster SUTVA and observation consistency replace the realized pattern and projected response in Definition \texttt{def:identified-response-laws} by their arm-\(a\) potentials on \(\{A_i=a\}\).  Hence
\[
\begin{aligned}
H_a(c,m,z_m)
 &=\nu^{-1}\mathbb E[V_i]\\
 &=\mathbb P^{\star}\{C_i=c,M_{ij}(a)=m,
      \operatorname{pr}_mY_{ij}(a)=z_m\}\\
 &=\sum_{s:\operatorname{pr}_m(s)=z_m}\lambda_a(c,s,m).
\end{aligned}
\]
Summing this identity over \(z_m\) and then over \(m\) gives
\[
\sum_{z_m}H_a(c,m,z_m)
 =\mathbb P^{\star}\{C_i=c,M_{ij}(a)=m\},
\qquad
\sum_{m,z_m}H_a(c,m,z_m)=\kappa(c).
\]
Because \(c\in\mathcal C_+\) means \(\kappa(c)>0\), Definition \texttt{def:identified-response-laws} therefore yields the well-defined equality
\[
r_a(m\mid c)=\mathbb P^{\star}\{M_{ij}(a)=m\mid C_i=c\}.
\]

If \(q_a(c,s)>0\), finite conditional-probability algebra gives
\[
\lambda_a(c,s,m)
=q_a(c,s)\mathbb P^{\star}\{M_{ij}(a)=m\mid C_i=c,Y_{ij}(a)=s\}.
\]
Multiplying the two inequalities in the pattern-selection sensitivity assumption by the positive number \(q_a(c,s)\) gives
\[
\Gamma^{-1}r_a(m\mid c)q_a(c,s)
\leq\lambda_a(c,s,m)
\leq\Gamma r_a(m\mid c)q_a(c,s).
\]
If \(q_a(c,s)=0\), nonnegativity and \(q_a(c,s)=\sum_m\lambda_a(c,s,m)\) force every \(\lambda_a(c,s,m)=0\), so the same inequalities hold without division.  Thus the tables form a point \(x\in\mathcal F_{\Gamma}(H)\).  Finally,
\[
\beta
=\sum_{c\in\mathcal C_+}\sum_{(s_0,s_1)\in R}\pi(c,s_0,s_1)
=c_R^{\mathsf T}x.
\]
For \(c\notin\mathcal C_+\), \(\kappa(c)=\mathbb P^{\star}(C_i=c)=0\), so the stratum has zero target mass.  We have proved the outer bound
\[
\Theta_I^{\Gamma}(H)\subseteq[L_{\Gamma},U_{\Gamma}].
\]

\emph{Technique: constructive reverse embedding of a feasible finite table.}
Fix an arbitrary \(x=(\lambda_0,\lambda_1,q_0,q_1,\pi)\in\mathcal F_{\Gamma}(H)\).  Summing the projection and completion equations, and using the observed-law calculation above, gives for every \(c\in\mathcal C_+\)
\[
\sum_s q_a(c,s)
=\sum_{s,m}\lambda_a(c,s,m)
=\sum_{m,z_m}H_a(c,m,z_m)
=\kappa(c).
\]
The coupling rows then give \(\sum_{s_0,s_1}\pi(c,s_0,s_1)=\kappa(c)\).  Define
\[
p_c(s_0,s_1)=\frac{\pi(c,s_0,s_1)}{\kappa(c)}
\]
and, on every cell with \(q_a(c,s)>0\), define
\[
\rho_a(m\mid c,s)=\frac{\lambda_a(c,s,m)}{q_a(c,s)}.
\]
The displayed positive denominators justify both divisions.  The coupling rows show that \(p_c\) is a probability mass function with arm-\(a\) marginal \(q_a(c,\cdot)/\kappa(c)\), and the completion row shows that each \(\rho_a(\cdot\mid c,s)\) is a probability mass function.  On a cell with \(q_a(c,s)=0\), choose any probability mass function for \(\rho_a(\cdot\mid c,s)\); the cell will be null.

Construct one cluster as follows.  Draw \(B_i\), including \(N_i\) and the complete labeled roster, from its observed preassignment distribution.  Since \(N_i\geq1\),
\[
0=\nu\kappa(c)=\mathbb E[N_i\mathbf 1\{C_i=c\}]
\quad\Longrightarrow\quad
\mathbb P(C_i=c)=0
\]
whenever \(c\notin\mathcal C_+\).  Thus \(C_i\in\mathcal C_+\) almost surely.  Conditional on a realized \(B_i\) with \(C_i=c\), for each roster label \(j\) draw
\[
(Y_{ij}(0),Y_{ij}(1))\sim p_c,
\]
and then, for each \(a\in\{0,1\}\), draw \(M_{ij}(a)\) conditionally from \(\rho_a(\cdot\mid c,Y_{ij}(a))\).  Conditional independence across labels and between the two response-pattern draws is one permissible construction; no maintained assumption requires such independence.  Finally draw \(A_i\) conditionally on \(B_i\) with probabilities \(e_a(B_i)\), independently of all potential quantities given \(B_i\), and form the realized record by observation consistency.  Repeating this construction independently across clusters yields the required i.i.d. cluster superpopulation, preserves the entire distribution of \(B_i\), satisfies cluster SUTVA and the known randomized assignment, and retains assignment positivity.

For any roster-member quantity \(f_{ij}\), the definition of individual weighting is
\[
\mathbb E^{\star}[f_{ij}]
=\nu^{-1}\mathbb E\!\left[\sum_{j=1}^{N_i}f_{ij}\right].
\]
Every label in stratum \(c\) has conditional cross-world law \(p_c\); consequently
\[
\begin{aligned}
&\mathbb P^{\star}\{C_i=c,Y_{ij}(0)=s_0,Y_{ij}(1)=s_1\}\\
&\qquad=\nu^{-1}\mathbb E[N_i\mathbf 1\{C_i=c\}]p_c(s_0,s_1)
=\kappa(c)p_c(s_0,s_1)
=\pi(c,s_0,s_1).
\end{aligned}
\]
Using the arm marginal of \(p_c\) and then the response draw gives
\[
\begin{aligned}
&\mathbb P^{\star}\{C_i=c,Y_{ij}(a)=s,M_{ij}(a)=m\}\\
&\qquad=\kappa(c)\frac{q_a(c,s)}{\kappa(c)}\rho_a(m\mid c,s)
=\lambda_a(c,s,m),
\end{aligned}
\]
where the final equality is immediate on positive \(q_a(c,s)\) cells and both sides are zero on null cells.  The projection equations therefore reproduce \(H_a\).  On every positive \(q_a(c,s)\) cell, division of the two feasible-set sensitivity inequalities gives
\[
\Gamma^{-1}r_a(m\mid c)
\leq\rho_a(m\mid c,s)
\leq\Gamma r_a(m\mid c),
\]
which is exactly the maintained selection restriction; that restriction imposes no condition after conditioning on a zero-probability state.  The constructed causal probability is \(c_R^{\mathsf T}x\).  Thus every feasible objective value is attainable by a compatible cluster law.

\emph{Technique: compact convex image.}
For each arm, summing all projection rows gives \(\sum_{c,s,m}\lambda_a(c,s,m)=1\); the completion and coupling rows then give \(\sum_{c,s}q_a(c,s)=1\) and \(\sum_{c,s_0,s_1}\pi(c,s_0,s_1)=1\).  Nonnegativity therefore puts every coordinate of \(x\) in \([0,1]\).  At fixed finite \(\Gamma\), the feasible set is the intersection of finitely many closed linear equalities and inequalities, so it is closed and bounded in a finite-dimensional space, hence compact.  It is convex because all constraints are linear.  Therefore \(c_R^{\mathsf T}x\) attains its minimum and maximum at some \(x_L,x_U\).  For \(t\in[0,1]\),
\[
x_t=(1-t)x_L+t x_U\in\mathcal F_{\Gamma}(H),
\qquad
c_R^{\mathsf T}x_t=(1-t)L_{\Gamma}+tU_{\Gamma}.
\]
The reverse embedding applied to \(x_t\) attains every point between the endpoints.  Combining this with the outer bound proves the sharp equality
\[
\Theta_I^{\Gamma}(H)=[L_{\Gamma},U_{\Gamma}]=\mathcal I_{\Gamma}(H).
\]

\emph{Technique: finite linear elimination and dual certificates.}
First, if \(\alpha\) is unrestricted and \(\eta\geq0\) satisfy
\(E^{\mathsf T}\alpha-G_{\Gamma}^{\mathsf T}\eta\leq c_R\), then every feasible \(x\) obeys
\[
\bar h^{\mathsf T}\alpha
=x^{\mathsf T}E^{\mathsf T}\alpha
\leq c_R^{\mathsf T}x+x^{\mathsf T}G_{\Gamma}^{\mathsf T}\eta
\leq c_R^{\mathsf T}x,
\]
where the last step uses \(G_{\Gamma}x\leq0\) and \(\eta\geq0\).  This is weak duality with the asserted signs.

For completeness, exactness and attainment follow by an elementary finite-elimination argument.  Split the equality \(Ex=\bar h\) into its two inequality orientations and project the finite system
\[
Ex=\bar h,
\qquad G_{\Gamma}x\leq0,
\qquad -x\leq0,
\qquad c_R^{\mathsf T}x\leq v
\]
onto the scalar \(v\).  At a single elimination step, inequalities with a positive coefficient on the eliminated coordinate give upper bounds, inequalities with a negative coefficient give lower bounds, and inequalities with zero coefficient are retained.  Existence of that coordinate is equivalent to every lower bound being no larger than every upper bound, together with the retained inequalities.  Each resulting comparison is a nonnegative linear combination of two old inequalities after division by positive coefficient magnitudes.  Induction over the finitely many coordinates therefore gives an exact finite inequality description of the projection, and every projected inequality is a nonnegative combination of the original inequality rows, with unrestricted multipliers for equality rows.

Compactness and nonemptiness show that this projection is exactly \([L_{\Gamma},\infty)\).  Consequently one of its finitely many lower inequalities binds at \(v=L_{\Gamma}\).  In the nonnegative combination producing that inequality, the coefficient on the objective row is positive; normalize it to one.  If \(\alpha\) is the resulting unrestricted equality multiplier, \(\eta\geq0\) is the multiplier on \(G_{\Gamma}x\leq0\), and \(\zeta\geq0\) is the multiplier on \(-x\leq0\), cancellation of the eliminated \(x\)-coordinates reads
\[
-E^{\mathsf T}\alpha+G_{\Gamma}^{\mathsf T}\eta-\zeta+c_R=0.
\]
Thus
\[
E^{\mathsf T}\alpha-G_{\Gamma}^{\mathsf T}\eta
=c_R-\zeta\leq c_R,
\qquad
\bar h^{\mathsf T}\alpha=L_{\Gamma}.
\]
This produces an attained dual optimizer and proves
\[
L_{\Gamma}
=\max_{\alpha\ \mathrm{unrestricted},\ \eta\geq0}
\left\{\bar h^{\mathsf T}\alpha:
E^{\mathsf T}\alpha-G_{\Gamma}^{\mathsf T}\eta\leq c_R\right\}.
\]
Applying the same argument to minimization of \(-c_R^{\mathsf T}x\), and replacing its unrestricted equality multiplier by its negative, yields the attained certificate
\[
U_{\Gamma}
=\min_{\alpha\ \mathrm{unrestricted},\ \eta\geq0}
\left\{\bar h^{\mathsf T}\alpha:
E^{\mathsf T}\alpha+G_{\Gamma}^{\mathsf T}\eta\geq c_R\right\}.
\]
This proves every assertion of the theorem.
\end{proof}
\par\smallskip\noindent \textit{Justification.} The theorem prevents loss of sharpness at the interface between multivariate missing-state completion and same-person causal coupling. \textit{Gap.} Scharfstein, Manski, and Anthony constrain scalar two-occasion follow-up laws and compare arm distributions. Gao, Ge, and Qian, Ji, Lei, and Spector, and Ober-Reynolds optimize over fixed or already available arm marginals. Fan, Park, Pass, and Shi characterize conditional optimal transport when the two source laws of \((Y_1,X)\) and \((Y_0,X)\) are identified. That two-source model and the present overlapping response-pattern model are non-nested, and none of these results characterizes the present latent arm-state completion polytope jointly with the cross-world Pareto coupling. \textit{Consumer.} A PPACT companion analysis can report a certified range for the fraction of enrolled patients whose complete treated outcome vector strictly Pareto-dominates their complete control vector despite incomplete component records.

\begin{proposition}[prop:nonrectangular-witness]\label{prop:nonrectangular-witness}
For the law \(\mathsf W\) at \(\Gamma=1\), the exact joint program gives \(\mathcal I_1(\mathsf W)=[1/6,8/9]\). The statewise-envelope rectangular relaxation gives upper endpoint \(1\): it admits the control vector \((1,1,1,0,2,0,1,0)/6\) and treatment vector \((0,0,0,0,1,1,2,2)/6\), which possess an all-benefit coupling, although the displayed control vector violates the overlapping pair-projection equations. This finite construction is an internal nonrectangularity diagnostic: replacing the exact completion polytope by its statewise rectangular envelope before transport loses response-pattern restrictions and strictly enlarges the feasible range.
\end{proposition}
\begin{proof}
There is one positive-mass response stratum, whose mass is one; suppress its index.  Order the binary states as
\[
000,001,010,011,100,101,110,111.
\]
We use finite-table elimination for the arm completions and elementary transport-cut inequalities for the coupling.  All displayed masses in the definition of \(\mathsf W\) are conditional on one of the three active response patterns, each of which has probability \(1/3\).  At \(\Gamma=1\), the two sensitivity inequalities in \(\mathcal F_1(\mathsf W)\) coincide.  Hence, for each active pattern \(m\),
\[
\lambda_a(s,m)=\frac{1}{3}q_a(s),
\]
whereas a pattern having \(r_a(m)=0\) has \(\lambda_a(s,m)=0\).  Substitution in the projection equations shows that the corresponding pairwise marginals of \(q_a\) must equal the conditional projection tables displayed in the definition of \(\mathsf W\).  Conversely, any nonnegative \(q_a\) with those pairwise marginals, together with \(\lambda_a(s,m)=q_a(s)/3\) on the three active patterns and zero elsewhere, satisfies every arm-specific constraint.  Thus the following elimination is an exact characterization, rather than a relaxation.

Write \(p_s=6q_0(s)\) and put \(t=p_{000}\).  Successively using the control margins for coordinate pairs \(12\), \(13\), and \(23\) gives
\[
\begin{aligned}
p_{001}&=1-t &&\text{from }p_{000}+p_{001}=1,\\
p_{010}&=1-t &&\text{from }p_{000}+p_{010}=1,\\
p_{100}&=2-t &&\text{from }p_{000}+p_{100}=2,\\
p_{011}&=t   &&\text{from }p_{010}+p_{011}=1,\\
p_{101}&=t   &&\text{from }p_{001}+p_{101}=1,\\
p_{110}&=1+t &&\text{from }p_{100}+p_{110}=3,\\
p_{111}&=1-t &&\text{from }p_{110}+p_{111}=2.
\end{aligned}
\]
The unused entries of all three pairwise tables are then satisfied by direct substitution.  Nonnegativity is equivalent to \(0\le t\le1\), and the entries sum to six.  Therefore every exact control completion, and no other one, is
\[
q_0^{,t}=\frac{1}{6}(t,1-t,1-t,t,2-t,t,1+t,1-t),
\qquad 0\le t\le1.
\]
For treatment, the zero cells in the \(12\) and \(13\) margins first force the four entries at \(000,001,010,011\) to vanish.  The \(23\) cells \(00\) and \(01\), the \(13\) cell \(10\), and the \(12\) cell \(11\) then give, respectively, the remaining masses \(1,1,2,2\) in units of \(1/6\).  Hence the treatment completion is unique:
\[
q_1=\frac{1}{6}(0,0,0,0,1,1,2,2).
\]
Both arm masses equal one.  Consequently a coupling exists for every \(t\in[0,1]\), for example the product coupling, and every such coupling has the support and normalization required by \(\mathcal F_1(\mathsf W)\).

We next bound the strict Pareto mass of an arbitrary coupling \(\pi\) of \(q_0^{,t}\) and \(q_1\).  Because \(111\) strictly dominates every binary state except itself,
\[
\begin{aligned}
\pi(R)
&\ge \sum_{s_0\ne111}\pi(s_0,111)\\
&=q_1(111)-\pi(111,111)\\
&\ge q_1(111)-q_0^{,t}(111)
=\frac{2}{6}-\frac{1-t}{6}
=\frac{1+t}{6}
\ge\frac16.
\end{aligned}
\]
This is the lower transport-cut certificate.  At \(t=0\), define a coupling by the following nonzero entries, each of mass \(1/6\):
\[
001\mapsto110,
\quad010\mapsto101,
\quad100\mapsto100,
\quad100\mapsto111,
\quad110\mapsto110,
\quad111\mapsto111.
\]
Its row vector is \((0,1,1,0,2,0,1,1)/6=q_0^{,0}\), and its column vector is \((0,0,0,0,1,1,2,2)/6=q_1\).  Of its six mass units, only \(100\mapsto111\) belongs to \(R\).  With \(\lambda_a=q_a/3\) on the active patterns as above, this is a point of \(\mathcal F_1(\mathsf W)\), so the exact lower endpoint is \(1/6\).

For the upper endpoint, no pair beginning at \(111\) belongs to the strict relation \(R\).  Therefore
\[
\pi(R)\le1-q_0^{,t}(111)=\frac{5+t}{6}.
\]
There is a second cut.  Let \(A=\{011,101,110,111\}\).  A strict Pareto edge beginning in \(A\) can end only at \(111\), so the benefit mass originating in \(A\) is at most \(q_1(111)=2/6\).  The benefit mass originating outside \(A\) is at most the available row mass
\[
q_0^{,t}(A^{\mathsf c})
=\frac{t+(1-t)+(1-t)+(2-t)}6
=\frac{4-2t}{6}.
\]
It follows that
\[
\pi(R)\le\frac{6-2t}{6}=1-\frac{t}{3}.
\]
Combining the two cut certificates,
\[
\pi(R)\le
\min\!\left\{\frac{5+t}{6},\frac{6-2t}{6}\right\}
\le\frac89.
\]
Indeed, the first affine bound is at most \(8/9\) when \(t\le1/3\), and the second is at most \(8/9\) when \(t\ge1/3\).

At \(t=1/3\), the following are the nonzero entries of a coupling, now written with their masses:
\[
\begin{array}{c|c}
(s_0,s_1)&\pi(s_0,s_1)\\ \hline
(000,100)&1/18\\
(001,101)&1/9\\
(010,110)&1/9\\
(011,111)&1/18\\
(100,101)&1/18\\
(100,110)&2/9\\
(101,111)&1/18\\
(110,111)&2/9\\
(111,100)&1/9
\end{array}
\]
The row masses are
\[
\frac1{18}(1,2,2,1,5,1,4,2)=q_0^{,1/3},
\]
and the column masses are \((0,0,0,0,1,1,2,2)/6=q_1\).  Every displayed edge except \((111,100)\) lies in \(R\); hence its benefit mass is \(1-1/9=8/9\).  This proves that the upper certificate is attained.  Since \(\mathcal F_1(\mathsf W)\) is convex and the objective is linear, convex mixtures of the lower- and upper-attaining feasible points attain every intermediate value.  Thus the exact identified range is
\[
\mathcal I_1(\mathsf W)=[1/6,8/9].
\]

It remains to check the claimed strict enlargement under the rectangular relaxation.  The exact parameterization gives the coordinatewise control envelopes
\[
6\ell_0=(0,0,0,0,1,0,1,0),
\qquad
6u_0=(1,1,1,1,2,1,2,1),
\]
while treatment is point identified, so \(\ell_1=u_1=q_1\).  The proposed relaxed control vector
\[
\widetilde q_0=\frac16(1,1,1,0,2,0,1,0)
\]
lies coordinatewise between \(\ell_0\) and \(u_0\), is nonnegative, and sums to one.  Couple it to \(q_1\) using the six entries
\[
000\mapsto100,
\quad001\mapsto101,
\quad010\mapsto110,
\quad100\mapsto110,
\quad100\mapsto111,
\quad110\mapsto111,
\]
each with mass \(1/6\).  These entries have row marginal \(\widetilde q_0\) and column marginal \(q_1\), and every one is a strict Pareto edge.  This feasible point of \(\mathcal F^{\mathrm{rect}}_1(\mathsf W)\) has objective one; normalization makes one an upper bound for every coupling, so the relaxed upper endpoint equals one.

Finally, \(\widetilde q_0\) is not an exact control completion.  For response pattern \(110\), its projected mass at pair state \(00\) is
\[
\widetilde q_0(000)+\widetilde q_0(001)=\frac{2}{6},
\]
whereas the corresponding control table in \(\mathsf W\) requires \(1/6\).  Equivalently, with \(\lambda_0=\widetilde q_0/3\), the left side of the exact projection equation would be \(1/9\), while its observed-law right side is \((1/3)(1/6)=1/18\).  Thus the relaxed witness violates an overlapping pair-projection equality, proving that the increase from \(8/9\) to one is caused precisely by discarding the joint cross-state projection restrictions.
\end{proof}
\par\smallskip\noindent \textit{Justification.} A strict finite gap establishes that the joint completion-and-coupling geometry is substantive rather than an artifact of notation. \textit{Gap.} The diagnostic makes no claim about a published imputation workflow. It isolates the information discarded by the internally defined statewise-envelope relaxation and thereby explains the theorem's nonrectangular feasible set. \textit{Consumer.} Analysts can use the finite table as a reproducible unit test for implementations of the exact joint program and its deliberately coarser rectangular benchmark.

\begin{proposition}[prop:majority-sensitivity-threshold]\label{prop:majority-sensitivity-threshold}
Suppose \(\Gamma_{\mathrm{feas}}<\infty\). On the feasible domain \(\Gamma\in[\Gamma_{\mathrm{feas}},\infty)\), the sets \(\mathcal F_{\Gamma}(H)\) are nested and \(L_{\Gamma}\) is nonincreasing. If \(\Gamma_{\mathrm{maj}}<\infty\), the infimum in its definition is attained. Guaranteed strict majority benefit holds exactly for feasible \(\Gamma<\Gamma_{\mathrm{maj}}\), equivalently \(L_{\Gamma}>1/2\); at and above a finite \(\Gamma_{\mathrm{maj}}\), an observationally compatible law has \(\beta\le1/2\). If \(\Gamma_{\mathrm{maj}}=\infty\), every finite feasible sensitivity value guarantees \(\beta>1/2\). Values \(\Gamma<\Gamma_{\mathrm{feas}}\) are model-data incompatibilities and carry no majority-benefit interpretation. If the entries of \(H_a\) are rational, each finite threshold is computable to any prescribed positive rational tolerance by monotone rational linear-feasibility searches; when a threshold is infinite, the same search supplies valid increasing lower bounds but need not terminate with a certificate of infinity.
\end{proposition}
\begin{proof}
Assume \(\Gamma_{\mathrm{feas}}<\infty\).  The coordinatewise comparison in the proof of Proposition \(\mathrm{prop:feasibility\mbox{-}onset}\) gives
\[
\mathcal F_{\gamma}(H)\subseteq\mathcal F_{\gamma'}(H)
\quad\text{whenever}\quad
\Gamma_{\mathrm{feas}}\le\gamma\le\gamma'<\infty.
\]
Minimizing the same linear objective over nested sets therefore gives \(L_{\gamma'}\le L_{\gamma}\).

Let
\[
\mathscr B=\{\gamma\in[1,\infty):
\exists x\in\mathcal F_{\gamma}(H),\ c_R^{\mathsf T}x\le1/2\}.
\]
This set is upward closed by the same nesting argument, and its infimum is \(\Gamma_{\mathrm{maj}}\) by definition.  If \(\mathscr B\ne\varnothing\), repeat the compactness-and-continuity argument used for the feasibility onset, now retaining the additional closed inequality \(c_R^{\mathsf T}x\le1/2\).  It follows that \(\mathscr B=[\Gamma_{\mathrm{maj}},\infty)\), so a finite majority threshold is attained.

For every feasible finite \(\gamma\), compactness and the sharpness theorem give
\[
\gamma\in\mathscr B
\quad\Longleftrightarrow\quad
L_{\gamma}\le\frac12
\quad\Longleftrightarrow\quad
\text{some compatible law has }\beta\le\frac12.
\]
If \(\gamma<\Gamma_{\mathrm{maj}}\) is feasible, membership in \(\mathscr B\) fails.  Since the minimum is attained, this is equivalent to \(L_{\gamma}>1/2\), and sharpness then forces \(\beta>1/2\) under every compatible law.  At and above a finite threshold, membership in \(\mathscr B\) and the reverse embedding from the sharpness theorem supply a compatible law with \(\beta\le1/2\).  If \(\Gamma_{\mathrm{maj}}=\infty\), \(\mathscr B\) is empty, so the strict inequality holds at every finite feasible index.  Proposition \(\mathrm{prop:feasibility\mbox{-}onset}\) separately shows that indices below \(\Gamma_{\mathrm{feas}}\) are incompatibilities, not majority conclusions.

Finally suppose that every entry of \(H_a\), and hence every \(r_a(m\mid c)\), is rational, and let \(\delta>0\) be rational.  At any positive rational trial value \(\gamma\), both predicates
\[
\mathcal F_{\gamma}(H)\ne\varnothing
\quad\text{and}\quad
\exists x\in\mathcal F_{\gamma}(H):c_R^{\mathsf T}x\le1/2
\]
are finite rational linear-feasibility problems and are therefore exactly decidable by finite elimination or a rational-pivot linear-program algorithm.  Starting at one and doubling the trial value finds an upper bracket whenever the corresponding threshold is finite; monotonicity then permits bisection until the rational bracket has width at most \(\delta\).  If a threshold is infinite, the successive failed trials remain valid increasing lower bounds, although this search need not terminate with a certificate of infinity.  This proves the stated, explicitly qualified computability claim.
\end{proof}
\par\smallskip\noindent \textit{Justification.} The threshold converts an indexed missingness departure into the exact point where a clinically interpretable majority claim ceases to be forced by the data. \textit{Gap.} Huang, Fang, Hanley, and Rosenblum study confidence intervals for a scalar ordinal or binary fraction benefiting in randomized trials, while existing cluster-trial sensitivity analyses target smooth average effects. Neither result propagates vector response-pattern selection through the same-person strict-Pareto event. \textit{Consumer.} A data-monitoring committee can state how large pattern-specific outcome selection must be before the conclusion that most patients benefit is no longer guaranteed.

\begin{proposition}[prop:complete-response-reduction]\label{prop:complete-response-reduction}
If \(M_{ij}(a)=(1,\ldots,1)\) almost surely in both arms, then \(\Gamma_{\mathrm{feas}}=1\), \(q_a=H_a\) is point identified, the sensitivity inequalities are redundant for every finite \(\Gamma\ge1\), and \(\mathcal I_{\Gamma}(H)\) reduces to the ordinary sharp optimal-transport interval for the strict Pareto cost with fixed arm marginals.
\end{proposition}
\begin{proof}
Let \(m_{\mathrm{all}}=(1,\ldots,1)\).  Under complete response, the definition of \(H_a\) gives
\[
H_a(c,m,z_m)=0\quad(m\ne m_{\mathrm{all}}),
\qquad
H_a(c,m_{\mathrm{all}},s)
=\mathbb P^{\star}\{C_i=c,Y_{ij}(a)=s\}.
\]
Consequently \(r_a(m_{\mathrm{all}}\mid c)=1\) and \(r_a(m\mid c)=0\) for every other pattern.  For \(m_{\mathrm{all}}\), the projection \(\operatorname{pr}_{m_{\mathrm{all}}}\) is the identity.  Its projection row therefore forces
\[
\lambda_a(c,s,m_{\mathrm{all}})=H_a(c,m_{\mathrm{all}},s).
\]
Every other projection right-hand side is zero, so nonnegativity forces \(\lambda_a(c,s,m)=0\) when \(m\ne m_{\mathrm{all}}\).  The completion equation now yields the point identification
\[
q_a(c,s)=\sum_m\lambda_a(c,s,m)
=H_a(c,m_{\mathrm{all}},s),
\]
which is the meaning of \(q_a=H_a\) in the complete-response reduction.

At \(\Gamma=1\), the sensitivity inequalities for the full pattern reduce to \(q_a\le\lambda_a\le q_a\), and those for every zero-frequency pattern reduce to \(0\le\lambda_a\le0\); all hold by the preceding identities.  For any finite \(\Gamma\ge1\), the full-pattern inequalities are
\[
\Gamma^{-1}q_a(c,s)\le q_a(c,s)\le\Gamma q_a(c,s),
\]
and all other pattern inequalities again say that zero equals zero.  Hence they impose no additional restriction.

For each \(c\in\mathcal C_+\), both fixed arm marginals have mass
\[
\sum_s q_a(c,s)=\sum_{m,z_m}H_a(c,m,z_m)=\kappa(c).
\]
Therefore a stratum-preserving coupling exists; for example,
\[
\pi(c,s_0,s_1)=
\frac{q_0(c,s_0)q_1(c,s_1)}{\kappa(c)}
\]
is well defined because \(\kappa(c)>0\).  Thus \(\mathcal F_1(H)\ne\varnothing\), and since the sensitivity index is restricted to \([1,\infty)\), the feasibility onset equals one.  After eliminating the forced \(\lambda_a\) and \(q_a\), the only remaining variables are precisely the nonnegative couplings \(\pi\) with these fixed arm marginals.  Minimizing and maximizing the strict-Pareto cost \(\mathbf 1_R\) over that transport polytope is the ordinary sharp finite optimal-transport interval.  The joint sharpness theorem proves that this optimization interval, rather than merely an outer relaxation, equals \(\mathcal I_{\Gamma}(H)\) for every finite \(\Gamma\ge1\).
\end{proof}
\par\smallskip\noindent \textit{Justification.} The reduction is the required sanity check and fixes the relationship to established marginal optimal-transport bounds. \textit{Gap.} The known complete-response result is recovered as a boundary case; the contribution is the nonrectangular observed-data completion that disappears in this case. \textit{Consumer.} Implementers obtain an exact regression test against standard finite optimal-transport software when all outcome components are observed.

\begin{proposition}[prop:joint-payoff-linear-image-corollary]\label{prop:joint-payoff-linear-image-corollary}
Fix a finite \(\Gamma\geq\Gamma_{\mathrm{feas}}\), a positive integer \(K\), and real payoff arrays \(w_h:\mathcal C_+\times\mathcal S\times\mathcal S\to\mathbb R\), \(h=1,\ldots,K\).  Under the same information contract and maintained restrictions as Theorem \texttt{thm:joint-observed-data-sharpness}, the sharp joint identified set of the \(K\)-vector whose \(h\)-th coordinate is
\[
\sum_{c,s_0,s_1}w_h(c,s_0,s_1)
\mathbb P^{\star}\{C_i=c,Y_{ij}(0)=s_0,Y_{ij}(1)=s_1\}
\]
is exactly
\[
\left\{
\left(\sum_{c,s_0,s_1}w_h(c,s_0,s_1)\pi(c,s_0,s_1)\right)_{h=1}^{K}
:x\in\mathcal F_{\Gamma}(H)
\right\}.
\]
This set is a nonempty compact convex polytope, and every one of its points is attained by a compatible cluster law.  For every direction \(v\in\mathbb R^K\), let \(c_v\) have entry \(\sum_{h=1}^K v_h w_h(c,s_0,s_1)\) on the coordinate \(\pi(c,s_0,s_1)\) and zero on every other coordinate of \(x\).  Its directional support value and corresponding dual certificate are both attained and satisfy
\[
\max_{x\in\mathcal F_{\Gamma}(H)}c_v^{\mathsf T}x
=\min_{\alpha\ \mathrm{unrestricted},\ \eta\geq0}
\left\{\bar h^{\mathsf T}\alpha:
E^{\mathsf T}\alpha+G_{\Gamma}^{\mathsf T}\eta\geq c_v\right\}.
\]
For \(K=1\) and \(w_1(c,s_0,s_1)=\mathbf 1\{(s_0,s_1)\in R\}\), the joint set reduces to \(\mathcal I_{\Gamma}(H)=[L_{\Gamma},U_{\Gamma}]\).
\end{proposition}
\begin{proof}
Define the linear map
\[
T_w(x)=\left(\sum_{c,s_0,s_1}w_h(c,s_0,s_1)\pi(c,s_0,s_1)\right)_{h=1}^{K}.
\]
The outer-map part of Theorem \texttt{thm:joint-observed-data-sharpness} sends every compatible cluster law to an \(x\in\mathcal F_{\Gamma}(H)\) whose \(\pi\)-coordinates equal that law's individual-weighted joint table of \((C_i,Y_{ij}(0),Y_{ij}(1))\).  Finite summation therefore gives its payoff vector as \(T_w(x)\), proving that the joint identified set is contained in \(T_w(\mathcal F_{\Gamma}(H))\).  Conversely, the theorem's reverse embedding sends every \(x\in\mathcal F_{\Gamma}(H)\) to a compatible cluster law with exactly that \(\pi\)-table.  The same finite-sum identity gives payoff vector \(T_w(x)\), proving the reverse inclusion and pointwise attainment.

The theorem proves that \(\mathcal F_{\Gamma}(H)\) is nonempty, compact, and convex.  Linearity and continuity of \(T_w\) therefore make its image nonempty, compact, and convex.  To verify the polytope claim without adding a representation assumption, append coordinates \(t\in\mathbb R^K\) and the linear equations \(t=T_w(x)\) to the finite system defining \(\mathcal F_{\Gamma}(H)\).  Eliminating the finitely many coordinates of \(x\) by the exact lower-bound/upper-bound comparisons used in the theorem's duality proof produces finitely many linear inequalities in \(t\) whose solution set is exactly \(T_w(\mathcal F_{\Gamma}(H))\).  That set is bounded because it is compact; hence it is a bounded polyhedron, that is, a compact convex polytope.

For \(v\in\mathbb R^K\), direct expansion gives
\[
v^{\mathsf T}T_w(x)
=\sum_{c,s_0,s_1}\left(\sum_{h=1}^{K}v_hw_h(c,s_0,s_1)\right)
\pi(c,s_0,s_1)
=c_v^{\mathsf T}x.
\]
Compactness gives an attained primal maximum.  The theorem's finite-elimination dual derivation uses no special property of \(c_R\); substituting the finite vector \(c_v\) throughout gives weak duality
\[
c_v^{\mathsf T}x
\leq x^{\mathsf T}E^{\mathsf T}\alpha+x^{\mathsf T}G_{\Gamma}^{\mathsf T}\eta
\leq\bar h^{\mathsf T}\alpha
\]
whenever \(E^{\mathsf T}\alpha+G_{\Gamma}^{\mathsf T}\eta\geq c_v\), followed by equality and an attained dual optimizer from the same exact projection argument.  This proves the displayed support formula.  Finally, when \(K=1\) and \(w_1=\mathbf 1_R\), the definition of \(c_R\) gives \(T_w(x)=c_R^{\mathsf T}x\); Theorem \texttt{thm:joint-observed-data-sharpness} identifies its image as \([L_{\Gamma},U_{\Gamma}]=\mathcal I_{\Gamma}(H)\).
\end{proof}
\par\smallskip\noindent \textit{Justification.} This corollary upgrades scalar reporting to a jointly attainable region for any fixed finite menu of clinically prespecified cross-world payoffs. \textit{Gap.} It exposes a consequence of the present joint completion-and-coupling geometry and does not add a network representation, an extremality assertion, or an inference guarantee. \textit{Consumer.} A PPACT companion analysis can report compatible combinations of several payoff summaries and optimize any prespecified direction without treating coordinatewise bounds as jointly attainable.

\section*{Technical internal limitation (diagnostic only)}

The finite rectangular-relaxation witness is an internal diagnostic of the proposed geometry, not a claim that a named published estimator or imputation procedure implements that relaxation. Sharpness is relative to the identified individual-weighted response-pattern laws \(H_0,H_1\) on \(\mathcal C_{+}\), the complete roster distribution, and the stated sensitivity restriction. A \(\kappa\)-null stratum has zero individual-weighted contribution and is deliberately absent from every conditional-law and linear-program index set. Structural zeros within positive-mass strata can make \(\Gamma_{\mathrm{feas}}=\infty\), in which case no finite member of the indexed sensitivity model is compatible with the observed law. Higher-order within-cluster observed-outcome dependence is not retained in this information contract and could further contract the identified set. The proposal makes no root-\(G\) inference, bootstrap, efficiency, or PPACT precision claim.

\section{Honest scope}

The compatibility-onset proposition, central sharpness theorem, joint-payoff linear-image corollary, finite nonrectangularity diagnostic, majority threshold, and complete-response reduction are proved. Every conditional response law and feasible-set row is restricted to \(\mathcal C_{+}\); null strata contribute zero target mass. Majority benefit is interpreted only for finite \(\Gamma\geq\Gamma_{\mathrm{feas}}\); smaller values represent infeasible response models. The identification core makes no finite-relation max-flow, projected-incidence obstruction, general network-representation, endpoint-extremality, or inference claim. Empirical usefulness at \(106\) clusters and access to a fixed PPACT cohort remain outside this identification core.

\begin{thebibliography}{99}

  \bibitem{WangLiWang2026} Wang, Bingkai, Fan Li, and Rui Wang. 2026. Handling incomplete outcomes and covariates in cluster-randomized trials: doubly robust estimation, efficiency considerations, and sensitivity analysis. Biometrics 82(1), ujag030. \(\operatorname{doi}:10.1093/biomtc/ujag030\).

  \bibitem{ScharfsteinManskiAnthony2004} Scharfstein, Daniel O., Charles F. Manski, and James C. Anthony. 2004. On the Construction of Bounds in Prospective Studies with Missing Ordinal Outcomes: Application to the Good Behavior Game Trial. Biometrics 60(1):154--164. \(\operatorname{doi}:10.1111/j.0006-341X.2004.00158.x\).

  \bibitem{JiLeiSpector2023} Ji, Wenlong, Lihua Lei, and Asher Spector. 2023. Model-Agnostic Covariate-Assisted Inference on Partially Identified Causal Effects. arXiv:2310.08115, current public revision checked in 2026.

  \bibitem{OberReynolds2023} Ober-Reynolds, Daniel. 2023. Estimating Functionals of the Joint Distribution of Potential Outcomes with Optimal Transport. Public working paper, arXiv:2311.09435, draft dated December 2, 2023.

  \bibitem{GaoGeQian2025} Gao, Zijun, Shu Ge, and Jian Qian. 2025. Bridging Multiple Worlds: Multi-marginal Optimal Transport for Causal Partial-identification Problem. Proceedings of the 28th International Conference on Artificial Intelligence and Statistics, Proceedings of Machine Learning Research 258:721--729.

  \bibitem{FanParkPassShi2025} Fan, Yanqin, Hyeonseok Park, Brendan Pass, and Xuetao Shi. 2025. Partial Identification in Moment Models with Incomplete Data---A Conditional Optimal Transport Approach. arXiv:2503.16098v2, revised October 2, 2025.

  \bibitem{LiMiaoShpitserTchetgen2023} Li, Yilin, Wang Miao, Ilya Shpitser, and Eric J. Tchetgen Tchetgen. 2023. A Self-Censoring Model for Multivariate Nonignorable Nonmonotone Missing Data. Biometrics 79(4):3203--3214. \(\operatorname{doi}:10.1111/biom.13916\).

  \bibitem{Strassen1965} Strassen, Volker. 1965. The Existence of Probability Measures with Given Marginals. The Annals of Mathematical Statistics 36(2):423--439. \(\operatorname{doi}:10.1214/aoms/1177700153\).

  \bibitem{LuDingDasgupta2018} Lu, Jiannan, Peng Ding, and Tirthankar Dasgupta. 2018. Treatment Effects on Ordinal Outcomes: Causal Estimands and Sharp Bounds. Journal of Educational and Behavioral Statistics 43(5):540--567. \(\operatorname{doi}:10.3102/1076998618776435\).

  \bibitem{HuangFangHanleyRosenblum2019} Huang, Emily J., Ethan X. Fang, Daniel F. Hanley, and Michael Rosenblum. 2019. Constructing a Confidence Interval for the Fraction Who Benefit from Treatment, Using Randomized Trial Data. Biometrics 75(4):1228--1239. \(\operatorname{doi}:10.1111/biom.13101\).

  \bibitem{WangParkSmallLi2024} Wang, Bingkai, Chan Park, Dylan S. Small, and Fan Li. 2024. Model-Robust and Efficient Covariate Adjustment for Cluster-Randomized Experiments. Journal of the American Statistical Association 119(548):2894--2907. \(\operatorname{doi}:10.1080/01621459.2023.2289693\).

  \bibitem{RabideauWang2021} Rabideau, Dustin J., and Rui Wang. 2021. Randomization-Based Confidence Intervals for Cluster Randomized Trials. Biostatistics 22(4):920--935. \(\operatorname{doi}:10.1093/biostatistics/kxaa007\).

  \bibitem{HayesMoulton2017} Hayes, Richard J., and Lawrence H. Moulton. 2017. Cluster Randomised Trials, second edition. Boca Raton: CRC Press.

  \bibitem{Rubin1974} Rubin, Donald B. 1974. Estimating Causal Effects of Treatments in Randomized and Nonrandomized Studies. Journal of Educational Psychology 66(5):688--701.

\end{thebibliography}

\end{document}